\scanpage{264}
\mspage{76}
% archive sheet number: 264
La \unsure{restr.} du carré \deleted{\unsure{envisagé}} substantiel cherché $(C)$ à la sous-catégorie $\deleted{\unsure{I'}}I_1$ est donc un diagramme substantiel de squelette

\diagram{$(\mathrm{Sq}\qquad$\begin{tikzcd}[ampersand replacement=\&,row sep=large,column sep=large]
e\arrow[d]\&\\
M\&e\arrow[l]
\end{tikzcd}\,,}

dont il se déduit (à isom can. près) en appliquant $i_*$ à ce dernier, de façon à le compléter en

\diagram{\begin{tikzcd}[ampersand replacement=\&,row sep=small,column sep=small]
e\arrow[d]\&\Omega M\arrow[l,dashed]\arrow[d,dashed]\\
M\&e\arrow[l]
\end{tikzcd}\,.}

Il reste donc à décrire une \unsure{substantiation|substantification} canonique du squelette (Sq), en termes de $M$. Cela nous conduit à regarder la sous-catégorie ouverte

\diagram{$I_0=\{0\}: \qquad$\begin{tikzcd}[ampersand replacement=\&,row sep=small,column sep=large]
{}\arrow[r,dotted]\&{}\\
0\,\bullet\arrow[u,dotted]\arrow[r,dotted]\&{}\arrow[u,dotted]
\end{tikzcd}}

de $I_1$ (donc aussi de $I$) réduite à l'objet initial, et l'inclusion

\diagram{\begin{tikzcd}[ampersand replacement=\&] I_0\arrow[r,hook,"i_0"]\&I_1\end{tikzcd}\,,}

et de \deleted{\unsure{déf}} substantifier le squelette Sq par $(i_0)_!(M)$.

On peut dire à \unsure{postériori|post'eriori} que, si on \unsure{connaît}\deleted{\unsure{ssait}} le carré substantiel, il n'y a \inserted{à isom. unique près} qu'une seule façon de l'obtenir comme un $i_{1*}i_{0!}(?)$, avec $?$ dans $\mathbb D(I_0)\simeq\mathcal A$, car $i_{0!}$ et $i_{1*}$ étant pl. fid., il en est de $=$ de leur composé. \deleted{Ainsi le \unsure{fait} qu'il}
